\documentclass[a4paper, serif, 10pt]{moderncv}

%% moderncv
% color and style
\moderncvcolor{black}
\moderncvstyle{banking}

%% page layout/margins
\usepackage[top=2.5cm, bottom=2.5cm, left=1.5cm, right=1.5cm]{geometry}

\nopagenumbers{}

%% Babel config for German language
% ref:
% - https://latex3.github.io/babel/guides/locale-german.html
\usepackage[ngerman]{babel}

%% fontspec
\usepackage{fontspec}

\IfFontExistsTF{Linux Libertine}{
  \setmainfont[Ligatures={TeX, Common}, Numbers=OldStyle]{Linux Libertine}
}{}
\IfFontExistsTF{Linux Libertine O}{
  \setmainfont[Ligatures={TeX, Common}, Numbers=OldStyle]{Linux Libertine O}
}{}

%% CTeX
% CJK support, used to show CJK characters in the resume
%
% - heading=false: disable ctex heading
% - punct=kaiming: use kaiming punctuations styles for CJK
% - scheme=plain: use plain scheme, do not override \normalsize font size
% - space=auto: space settings for CJK characters
%
% CJK fonts are resolved per locale language, e.g. Noto Serif/Sans CJK JP
% for Japanese resumes, with fallback to the next region if not installed
%
% ref:
% - http://ctan.mirrorcatalogs.com/language/chinese/ctex/ctex.pdf
\usepackage[UTF8, heading=false, punct=kaiming, scheme=plain, space=auto]{ctex}

\IfFontExistsTF{Noto Serif CJK SC}{
  \setCJKmainfont{Noto Serif CJK SC}
}{}
\IfFontExistsTF{Noto Sans CJK SC}{
  \setCJKsansfont{Noto Sans CJK SC}
}{}

%% line spacing
\usepackage{setspace}
\setstretch{1.125}

%% URL styling - use normal text instead of monospace
\urlstyle{same}

%% Patch \httplink and \httpslink to handle full URLs
% The original moderncv macros always prepend http:// or https:// to URLs.
% This causes issues when the URL already has a protocol (e.g., https://example.com)
% resulting in malformed URLs like https://https://example.com.
% This patch checks if the URL already contains "://" and uses it directly if so.
% Note: We use \str_set:Nx (x-expansion) to expand macros like \@homepage first.
\ExplSyntaxOn
\str_new:N \g_yamlresume_url_str

\RenewDocumentCommand{\httpslink}{O{}m}{
  \str_set:Nx \g_yamlresume_url_str {#2}
  \str_if_in:NnTF \g_yamlresume_url_str {://}
    { \url{#2} }
    { \url{https://#2} }
}

\RenewDocumentCommand{\httplink}{O{}m}{
  \str_set:Nx \g_yamlresume_url_str {#2}
  \str_if_in:NnTF \g_yamlresume_url_str {://}
    { \url{#2} }
    { \url{http://#2} }
}
\ExplSyntaxOff

\name{Lukas Hoffmann}{}
\title{Senior Data Engineer – Cloud \& Big Data}
\phone[mobile]{+49 30 12345678}
\email{lukas.hoffmann@example.de}
\homepage{https://www.lukashoffmann.de}

\address{Friedrichstraße 123, Berlin, Berlin, Deutschland, 10117}{}{}

\extrainfo{{\small \faGithub}\ \href{https://github.com/lukashoffmann}{@lukashoffmann} {} {} {} • {} {} {} 
{\small \faLinkedin}\ \href{https://www.linkedin.com/in/lukashoffmann}{@lukashoffmann}}

\begin{document}

\maketitle

\section{Basisinformationen}

\cvline{}{\begin{itemize}
\item Data Engineer mit über 6 Jahren Erfahrung in der Konzeption, Entwicklung und dem Betrieb skalierbarer Datenpipelines
\item Spezialisiert auf Cloud-Datenplattformen (AWS, GCP), ETL-Prozesse und die Modellierung von Data-Warehouse-Strukturen
\item Fundierte Kenntnisse in Python, SQL, Apache Spark, Kafka und Airflow sowie in der Automatisierung von Datenqualitätsprüfungen
\item Teamorientiert, analytisch denkend und leidenschaftlich für saubere, wartbare und performante Datenarchitekturen
\end{itemize}}
\section{Ausbildung}

\cventry{Okt. 2014–Sept. 2017}
        {Master, Informatik, Punkte: 1.3}
        {Technische Universität München}
        {\url{https://www.tum.de}}
        {}
        {\begin{itemize}
\item Schwerpunkt auf Datenbanken, verteilten Systemen und Big-Data-Technologien
\item Masterarbeit: „Entwurf einer skalierbaren Echtzeit-Datenpipeline für IoT-Anwendungen“
\item Praktische Erfahrung durch Projekte in Kooperation mit Industriepartnern
\end{itemize}
\textbf{Kurse}: Datenbanksysteme, Verteilte Systeme, Maschinelles Lernen, Big Data Analytics, Datenstromverarbeitung, Cloud Computing, Software Engineering}

\cventry{Okt. 2011–Sept. 2014}
        {Bachelor, Wirtschaftsinformatik, Punkte: 1.7}
        {Universität Hamburg}
        {\url{https://www.uni-hamburg.de}}
        {}
        {\begin{itemize}
\item Grundlagen in Programmierung, Datenbanken und betriebswirtschaftlicher Informatik
\item Bachelorarbeit: „Vergleich von ETL-Werkzeugen für heterogene Datenquellen“
\item Mitglied der Fachschaft Informatik und Organisation von Workshops zu SQL
\end{itemize}
\textbf{Kurse}: Programmierung, Datenbanken, Betriebswirtschaftslehre, Statistik, Informationssysteme}
\section{Berufserfahrung}

\cventry{Jan. 2020–Heute}
        {Senior Data Engineer}
        {DataCraft GmbH}
        {\url{https://www.datacraft.de}}
        {}
        {\begin{itemize}
\item Leitung eines Teams von vier Data Engineers bei der Migration einer monolithischen Datenplattform in eine Cloud-native Architektur auf AWS
\item Konzeption und Implementierung von ETL-Pipelines mit Apache Airflow, Spark und dbt zur Verarbeitung von über 5 TB täglich
\item Einführung automatisierter Datenqualitätsprüfungen und Monitoring-Dashboards zur Sicherstellung der Datenintegrität
\item Enge Zusammenarbeit mit Data Scientists, Analytics-Teams und Stakeholdern zur Ableitung datengetriebener Entscheidungen
\item Mentoring von Junior-Engineers und Förderung von Best Practices in den Bereichen Code-Qualität und Dokumentation
\end{itemize}
\textbf{Schlüsselwörter}: Datenpipelines, ETL, AWS, Apache Spark, Airflow, Data Warehouse}

\cventry{Okt. 2017–Dez. 2019}
        {Data Engineer}
        {CloudMetrics AG}
        {\url{https://www.cloudmetrics.de}}
        {}
        {\begin{itemize}
\item Entwicklung und Wartung von Batch- und Streaming-Pipelines mit Python, SQL und Apache Kafka
\item Aufbau eines Data Warehouse auf Google BigQuery inklusive Datenmodellierung und Performance-Optimierung
\item Automatisierung von Deployments mit Docker, Kubernetes und Terraform
\item Durchführung von Datenanalysen und Erstellung von Reports für das Management
\end{itemize}
\textbf{Schlüsselwörter}: Python, SQL, Kafka, BigQuery, GCP, Terraform}
\section{Sprachen}

\cvline{Deutsch}{Muttersprache oder zweisprachig \hfill \textbf{Schlüsselwörter}: Muttersprache}
\cvline{Englisch}{Verhandlungssichere Kenntnisse \hfill \textbf{Schlüsselwörter}: C1, Business fluent}
\section{Fähigkeiten}

\cvline{Programmierung \& Datenbanken}{Experte \hfill \textbf{Schlüsselwörter}: Python, SQL, Scala, PostgreSQL, MongoDB, Redis}
\cvline{Big Data \& Streaming}{Erfahren \hfill \textbf{Schlüsselwörter}: Apache Spark, Apache Kafka, Apache Flink, Hadoop, Pandas}
\cvline{Cloud \& DevOps}{Erfahren \hfill \textbf{Schlüsselwörter}: AWS, GCP, Docker, Kubernetes, Terraform, Airflow}
\cvline{Data Warehousing \& BI}{Erfahren \hfill \textbf{Schlüsselwörter}: Snowflake, BigQuery, Redshift, dbt, Tableau, Power BI}
\section{Auszeichnungen}

\cventry{Dez. 2023}
        {DataCraft GmbH}
        {Mitarbeiter des Jahres – Data Engineering}
        {}
        {}
        {Ausgezeichnet für die erfolgreiche Migration der Datenplattform in die AWS-Cloud und die Einführung automatisierter Datenqualitätsprüfungen.}
\section{Zertifikate}

\cventry{Mai 2022}
        {Amazon Web Services}
        {AWS Certified Data Analytics – Specialty}
        {\url{https://aws.amazon.com/certification/}}
        {}
        {}

\cventry{Aug. 2021}
        {Google Cloud}
        {Google Cloud Professional Data Engineer}
        {\url{https://cloud.google.com/certification/data-engineer}}
        {}
        {}
\section{Veröffentlichungen}

\cventry{Nov. 2021}
        {Skalierbare Echtzeit-Datenpipelines mit Kafka und Spark}
        {Datenbank-Spektrum}
        {\url{https://link.springer.com/journal/13222}}
        {}
        {Fachartikel über Architekturmuster und Betriebserfahrungen mit Streaming-Datenpipelines in der Cloud.}
\section{Projekte}

\cventry{März 2022–Juni 2023}
        {Aufbau einer skalierbaren Streaming-Plattform zur Verarbeitung von Sensordaten in Echtzeit.}
        {Echtzeit-Analyseplattform für IoT-Daten}
        {\url{https://github.com/lukashoffmann/iot-analytics}}
        {}
        {\begin{itemize}
\item Architektur und Implementierung einer Echtzeit-Pipeline mit Kafka, Spark Structured Streaming und AWS Kinesis
\item Verarbeitung von ca. 10.000 Ereignissen pro Sekunde mit Sub-Sekunden-Latenz
\item Aufbau eines Dashboards zur Visualisierung von Anomalien und Betriebskennzahlen
\end{itemize}
\textbf{Schlüsselwörter}: IoT, Kafka, Spark, AWS, Echtzeit}

\cventry{Jan. 2021–Heute}
        {Entwicklung eines Python-Frameworks zur automatisierten Datenqualitätsprüfung in ETL-Pipelines.}
        {Open-Source Data Quality Framework}
        {\url{https://github.com/lukashoffmann/dq-framework}}
        {}
        {\begin{itemize}
\item Open-Source-Bibliothek zur Definition und Ausführung von Datenqualitätsregeln in Python
\item Integration in Airflow und dbt zur automatischen Validierung von Pipeline-Ergebnissen
\item Über 200 Sterne auf GitHub und Nutzung durch mehrere Unternehmen
\end{itemize}
\textbf{Schlüsselwörter}: Python, Data Quality, Open Source, Airflow}
\section{Interessen}

\cvline{Open Source}{Apache Spark, Airflow, dbt}
\cvline{Sport}{Laufen, Radfahren, Wandern}
\section{Ehrenamtliche Tätigkeiten}

\cventry{Sept. 2020–Heute}
        {Dozent für Datenanalyse}
        {Data Literacy e.V.}
        {\url{https://www.data-literacy.de}}
        {}
        {\begin{itemize}
\item Durchführung von Workshops zu SQL und Python für Mitarbeitende von Non-Profit-Organisationen
\item Entwicklung von Lehrmaterialien und Übungsprojekten zur Datenanalyse
\item Beratung von Vereinen bei der Auswahl und Einrichtung von Datenbanken
\end{itemize}}
    
\end{document}